%% ma: L10-B14-0041
%% lop: 10
%% chuong: 5
%% bai: 14
%% dang: 10.14.7
%% loai: TL
%% muc: TH
%% dap_an: "-"
%% nguon: "chuyen-de-cac-so-dac-trung-cua-mau-so-lieu-khong-ghep-nhom-toan-10-chuong-trinh-moi.pdf (GV. Phan Nhật Linh) (P05), Bài 03 Dạng 2 Bài tập tự luận 1"
%% kien_thuc: "Lớp 10 Bài 13 (số trung bình), Bài 14 (phương sai, độ lệch chuẩn)"
%% trang_thai: chinh-thuc
%% ghi_chu: "Tự duyệt theo yêu cầu giáo viên 22/10/2026."
\begin{ex}
Tiền lương của $6$ nhân viên trong một văn phòng như sau (đơn vị: triệu đồng):
\begin{center}
\begin{tabular}{|c|c|c|c|c|c|}\hline
$5{,}5$ & $8$ & $7{,}5$ & $17$ & $13$ & $9$\\\hline
\end{tabular}
\end{center}
Tính phương sai của mẫu số liệu.
\loigiai{$\overline x=\dfrac{60}6=10$.

\[s^2=\dfrac{(5{,}5-10)^2+(8-10)^2+(7{,}5-10)^2+(17-10)^2+(13-10)^2+(9-10)^2}6=\dfrac{179}{12}\approx14{,}92.\]}
\end{ex}
